\section{Gemini Decisions API、判断専用モデルを発表}
GoogleはGemini Decision 1を利用するGemini Decisions APIを発表した。Grok原文によると、テキスト生成を挟まず、共通の入力に対して最大1,000件の二択・選択・採点の問いを一度に評価し、確率や信頼度を返す。マルチモーダル入力に対応し、応答時間は150～250ミリ秒、入力料金は100万トークン当たり0.10ドル、出力課金はないとされる。エージェントの経路選択や安全チェックへの応用が注目される一方、実運用での独立比較はこれからだ。
\dailyxsource{Google公式の発表}{https://x.com/Google/status/2120583966272983331}

\section{Claude Fable／Mythosの制作デモが拡大}
AnthropicのClaude Fable 5.2とMythos 5.2について、ゲーム制作、BlenderやUnreal向けの3Dアセット、科学研究、コーディングのデモが相次いだ。数時間でFPSゲームを作成したとの投稿や、数学の形式化への言及もある。価格や速度の改善を評価する声とともに、使用制限や誤生成を指摘する声も見られた。モデルの公式発表は11月9日で観測窓外。本号が記録するのは、その後の利用報告とコミュニティの反応である。

\section{Gemini 4 Carbonの利用報告とモデル比較}
Gemini 4 Carbonが有料開発者やUltra利用者へ広く開放された後、コーディングや科学タスクでの使用感が共有された。Fable系モデルとの比較やベンチマーク順位の変動が議論される一方、医療分野の専門家探索など特定タスクでの限界を挙げる投稿もある。Grok原文では長時間タスク向けの設計や、入力100万トークン当たり2ドルの導入価格に触れているが、実務上の優劣はタスクと評価条件に依存する。

\section{OpenAIの四半期売上を巡る報道}
BloombergやWSJを基にした投稿では、OpenAIの2026年第3四半期売上が約100億ドルで前期比約50\%増、調整後営業損失は約109億ドルと報じられた。年換算売上規模として500億ドルという数字も流通した。企業向け需要の拡大やAnthropicとの比較に関心が集まるが、Grokの収集時点でOpenAIの公式確認はなく、匿名情報源や投資家向け資料に基づく報道として留保が必要である。

\section{AIを活用した天文学の発見候補と助成}
公開望遠鏡データをClaudeとCodexで解析し、約10年間見落とされていたとする矮惑星候補を見つけたとの報告が投稿された。あわせて、アマチュア天文学研究に10万ドルを提供する助成制度が発表された。天文学者からの肯定的な反応も紹介されているが、天体の発見は確認作業が進行中であり、公式な天文機関による認定はGrok原文では未確認である。AIによる公開データの再解析と人間の検証の組み合わせが注目点となる。
\dailyxsource{発見候補と助成を紹介した投稿}{https://x.com/willdepue/status/2120632732107907471}

\section{MythosによるAMD GPU移植の未確認情報}
Mythosを使ってClaudeの実行環境を3日間でAMD GPUへ移植し、ほぼすべてのカーネルをAIが記述したとの主張が共有された。事実であればGPU向けソフトウェアの移植作業やCUDA依存を考える材料になる。ただし、Grokはこの話題をUnverifiedと判定しており、AnthropicやAMDの公式説明、コードや性能測定の独立検証は確認できていない。成果として確定せず、投稿上の主張として記録する。
\dailyxsource{移植を巡る関連投稿}{https://x.com/kshot9000/status/2120637175498776756}

\section{AnthropicのIPOを巡る市場の言及}
AnthropicのIPOが近いとの観測や、応募超過を示唆する報道・市場関連投稿が見られた。企業評価額や投資家の需要を巡る議論が広がっているが、Grok原文では正式な上場日程や申込状況の公式確認は示されていない。IPOの進行を既定事実とせず、モデル開発企業の資金調達への関心を表す未確認情報として扱う。

\section{Claude Codeの無料プロキシとOSS代替}
Claude Codeの利用コストや制限を背景に、無料プロキシやオープンソースモデルへ切り替える方法を紹介する投稿が観測された。複数のモデル提供元を切り替えれば費用を抑えられるとの主張がある一方、代替先がClaudeそのものではない点や、品質・セキュリティの違いも指摘された。紹介された個別ツールの安全性や継続的な利用条件は未検証であり、利用時には別途確認が必要だ。

\section{ベンチマーク順位を巡る議論}
Artificial Analysisなどの評価指数を引用し、Fable 5.2が上位でGemini 4 Carbonが続き、OpenAIのモデルは相対的に後れを取るとの比較投稿が流れた。OpenAIが約10週間にわたり新しいフロンティアモデルを出していないとの指摘もある。いずれもGrokが収集した順位表や投稿者の見解であり、測定条件、評価時点、モデルの利用可能性を揃えた独立比較を本号で行ったわけではない。